\documentclass[12pt, a4paper]{article}

% 1. Cấu hình Tiếng Việt và Font Times New Roman 13pt chuẩn
\usepackage[utf8]{inputenc}
\usepackage[T5]{fontenc}
\usepackage[vietnamese]{babel}
\usepackage{newtxtext,newtxmath}

\makeatletter
\renewcommand{\normalsize}{\@setfontsize\normalsize{13pt}{16pt}}
\makeatother
\normalsize

% 2. Gói Toán học & Ký hiệu bổ trợ
\usepackage{amsmath, amssymb, amsfonts, mathrsfs, amsthm}
\usepackage{graphicx}
\usepackage{fontawesome}
\usepackage{booktabs, tabularx, array, multirow}
\usepackage{enumitem}

% 3. Đồ họa TikZ, Bảng biến thiên & Hình không gian
\usepackage{tikz}
\usetikzlibrary{calc, intersections, angles, quotes, patterns, arrows.meta, 3d}
\usepackage{pgfplots}
\pgfplotsset{compat=1.18}
\usepackage{tkz-euclide}
\usepackage{tkz-tab}

\renewcommand{\vec}[1]{\overrightarrow{#1}}

% 4. Hộp sư phạm (Tcolorbox)
\usepackage[many]{tcolorbox}
\tcbuselibrary{skins, breakable}

\newtcolorbox{pedabox}[1][]{
    enhanced,
    colback=blue!5!white,
    colframe=blue!75!black,
    fonttitle=\bfseries\sffamily,
    title=#1,
    boxrule=0.8pt,
    arc=2mm,
    breakable,
    left=3mm, right=3mm, top=2mm, bottom=2mm
}

\newtcolorbox{scaffoldbox}[1][]{
    enhanced,
    colback=orange!5!white,
    colframe=orange!85!black,
    fonttitle=\bfseries\sffamily,
    title=#1,
    boxrule=0.8pt,
    arc=2mm,
    breakable,
    left=3mm, right=3mm, top=2mm, bottom=2mm
}

\newtcolorbox{aibox}[1][]{
    enhanced,
    colback=green!5!white,
    colframe=teal!80!black,
    fonttitle=\bfseries\sffamily,
    title=#1,
    boxrule=0.8pt,
    arc=2mm,
    breakable,
    left=3mm, right=3mm, top=2mm, bottom=2mm
}

% 5. Căn lề chuẩn văn bản giáo án
\usepackage{geometry}
\geometry{
    a4paper,
    left=25mm,
    right=15mm,
    top=20mm,
    bottom=20mm
}

% 6. Header / Footer chuẩn hóa (BẮT BUỘC CHÍNH XÁC NỘI DUNG NÀY)
\usepackage{fancyhdr}
\usepackage{lastpage}
\pagestyle{fancy}
\fancyhf{}

\fancyhead[L]{\small \textit{Kế hoạch bài dạy Toán 11}}
\fancyhead[R]{\textbf{Trang \thepage/\pageref{LastPage}}}
\fancyfoot[L]{\small \textit{Tổ Toán - Tin}}
\fancyfoot[R]{\small \textit{Trường THCS\&THPT Hồng Vân}}
\renewcommand{\headrulewidth}{0.4pt}
\renewcommand{\footrulewidth}{0.4pt}

% 7. Giãn dòng & Giãn đoạn theo chuẩn Nghị định 30/2020/NĐ-CP
\usepackage{setspace}
\setstretch{1.15}
\setlength{\parindent}{1.0cm}
\setlength{\parskip}{5pt}

\begin{document}

\begin{center}
    {\Large \textbf{KẾ HOẠCH BÀI DẠY MÔN TOÁN LỚP 11}}\\[0.4em]
    {\LARGE \textbf{BÀI TẬP CUỐI CHƯƠNG II}}\\[0.3em]
    {\large \textbf{(DÃY SỐ. CẤP SỐ CỘNG VÀ CẤP SỐ NHÂN)}}\\[0.5em]
    \textbf{Môn học:} Toán 11 | \textbf{Thời lượng:} 01 tiết (Tiết 17 theo PPCT)
\end{center}

\vspace{0.5em}
\hrule
\vspace{1em}

\section*{I. MỤC TIÊU DẠY HỌC}

\subsection*{1. Kiến thức, Năng lực}

\textbf{a) Năng lực Toán học đặc thù (nguyên văn chuẩn CT GDPT 2018):}
\begin{itemize}[leftmargin=1.5em]
    \item \textbf{Năng lực tư duy và lập luận toán học:}
    \begin{itemize}
        \item Hệ thống hóa và củng cố toàn bộ các khái niệm cốt lõi của Chương II: Dãy số, cách cho dãy số (liệt kê, công thức số hạng tổng quát, hệ thức truy hồi, mô tả); dãy số tăng, dãy số giảm, dãy số bị chặn (chặn trên, chặn dưới, bị chặn).
        \item So sánh, phân biệt bản chất toán học giữa Cấp số cộng (tăng trưởng theo phép cộng với công sai $d$ không đổi --- hàm bậc nhất) và Cấp số nhân (tăng trưởng theo phép nhân với công bội $q$ không đổi --- hàm số mũ).
        \item Khắc sâu và giải thích được các mối quan hệ đặc trưng:
        \begin{itemize}
            \item Ba số liên tiếp $a, b, c$ lập thành cấp số cộng khi và chỉ khi: $a + c = 2b$.
            \item Ba số liên tiếp $a, b, c$ lập thành cấp số nhân khi và chỉ khi: $a \cdot c = b^2$.
        \end{itemize}
        \item Nhận diện và giải thích bản chất công thức tính tổng của $n$ số hạng đầu tiên: $S_n = \dfrac{n(u_1 + u_n)}{2}$ đối với cấp số cộng và $S_n = \dfrac{u_1(1 - q^n)}{1 - q}$ ($q \neq 1$) đối với cấp số nhân.
    \end{itemize}
    \item \textbf{Năng lực giải quyết vấn đề toán học:}
    \begin{itemize}
        \item Thành thạo kĩ năng giải hệ phương trình hai ẩn $u_1, d$ (đối với cấp số cộng) và $u_1, q$ (đối với cấp số nhân) thông qua các điều kiện ràng buộc giữa các số hạng.
        \item Phân tích và giải quyết các bài toán tổng hợp: bài toán đan xen giữa cấp số cộng và cấp số nhân; bài toán tìm số hạng, xác định vị trí của số hạng, tính tổng các dãy số quy về cấp số.
        \item Vận dụng các công thức cấp số để tính toán các đại lượng thực tế.
    \end{itemize}
    \item \textbf{Năng lực mô hình hóa toán học:}
    \begin{itemize}
        \item Mô hình hóa các tình huống thực tiễn thành bài toán cấp số: bài toán tiết kiệm tích lũy định kì, bài toán lãi kép ngân hàng, bài toán xếp chồng vật liệu dạng tháp, bài toán khấu hao tài sản hàng năm.
    \end{itemize}
    \item \textbf{Năng lực giao tiếp toán học:}
    \begin{itemize}
        \item Trình bày lời giải bài toán một cách chặt chẽ, mạch lạc; sử dụng chính xác các kí hiệu toán học $u_n, d, q, S_n, \sum$. Tích cực phản biện, trao đổi kết quả với bạn học trong hoạt động nhóm.
    \end{itemize}
    \item \textbf{Năng lực sử dụng công cụ, phương tiện học toán:}
    \begin{itemize}
        \item Khai thác tối ưu máy tính cầm tay (MTCT Casio fx-580VN X) để giải hệ phương trình bậc nhất hai ẩn (\texttt{MENU} 9-1-2), tính giá trị các lũy thừa lớn và tính tổng dãy số bằng chức năng $\sum$ (\texttt{SHIFT} $x$).
    \end{itemize}
\end{itemize}

\textbf{b) Năng lực số (theo Thông tư 02/2025/TT-BGDĐT):}
\begin{itemize}[leftmargin=1.5em]
    \item \textbf{Mã 2.2NC1a (Chia sẻ thông tin và dữ liệu số):} Học sinh tạo lập sơ đồ tư duy hoặc bảng tóm tắt công thức Chương II (CSC và CSN) bằng công cụ số, thực hiện chia sẻ tài liệu ôn tập qua nhóm học tập Zalo của lớp hoặc thư mục chia sẻ Google Drive của trường; rèn luyện văn hóa tương tác và chia sẻ tài nguyên học tập số an toàn, đúng quy định.
\end{itemize}

\textbf{c) Năng lực AI (theo Quyết định 2422/QĐ-BGDĐT):}
\begin{itemize}[leftmargin=1.5em]
    \item \textbf{Mã 11.C3.1 (Kỹ thuật câu lệnh và khai thác trợ lý AI có cấu trúc):} Học sinh thực hành kĩ thuật đặt câu lệnh (prompt) theo phương pháp "Phân rã vấn đề" (Problem Decomposition) cho trợ lý AI (Gemini, ChatGPT) để chia nhỏ bài toán cấp số thành các câu hỏi gợi ý bước một (Step-by-step scaffolding), giúp học sinh tự tư duy từng bước thiết lập hệ phương trình thay vì yêu cầu AI giải hộ trọn gói; qua đó hình thành thói quen kiểm soát và làm chủ công nghệ AI.
\end{itemize}

\textbf{d) Năng lực chung \& Phẩm chất:}
\begin{itemize}[leftmargin=1.5em]
    \item \textbf{Năng lực chung:} Năng lực tự chủ và tự học (tự hệ thống kiến thức toàn chương, tự rà soát lỗ hổng kiến thức cá nhân); năng lực giao tiếp và hợp tác (thảo luận cặp đôi giải quyết bài toán khó).
    \item \textbf{Phẩm chất:} Chăm chỉ, cẩn thận, tỉ mỉ khi biến đổi công thức; trung thực trong quá trình tự đánh giá và đối chiếu lời giải với trợ lý học tập.
\end{itemize}

\section*{II. THIẾT BỊ DẠY HỌC VÀ HỌC LIỆU}
\begin{itemize}[leftmargin=1.5em]
    \item \textbf{Giáo viên:} Kế hoạch bài dạy chi tiết chuẩn CV 5512; bài trình chiếu PowerPoint hệ thống hóa kiến thức Chương II; Phiếu học tập số 1 (Bảng đối chiếu song song CSC --- CSN), Phiếu học tập số 2 (Bộ bài tập luyện tập phân hóa); máy tính, máy chiếu kết nối mạng Internet.
    \item \textbf{Học sinh:} Vở ghi bài, SGK Toán 11; máy tính cầm tay Casio fx-580VN X; thiết bị thông minh (điện thoại hoặc máy tính bảng) phục vụ chia sẻ tài liệu số và thực hành câu lệnh AI.
\end{itemize}

\section*{III. TIẾN TRÌNH DẠY HỌC (01 TIẾT --- 45 PHÚT)}

\begin{center}
    \framebox[1.0\textwidth]{\parbox{0.96\textwidth}{
        \textbf{CẤU TRÚC TIẾN TRÌNH TIẾT 17 (TIẾT BÀI TẬP CUỐI CHƯƠNG II):}\\[0.3em]
        $\bullet$ \textbf{Hoạt động 1 (05 phút):} Khởi động --- Trò chơi "Phân loại thần tốc" nhận diện nhanh CSC và CSN.\\[0.3em]
        $\bullet$ \textbf{Hoạt động 2 (12 phút):} Ôn tập hệ thống hóa kiến thức \& Bảng đối chiếu CSC --- CSN --- Tích hợp Năng lực số (Mã 2.2NC1a) \& Năng lực AI (Mã 11.C3.1).\\[0.3em]
        $\bullet$ \textbf{Hoạt động 3 (20 phút):} Luyện tập giải toán phân hóa (3 dạng bài tập cốt lõi).\\[0.3em]
        $\bullet$ \textbf{Hoạt động 4 (08 phút):} Vận dụng thực tiễn --- Bài toán so sánh hai mô hình tích lũy tài chính.
    }}
\end{center}

\vspace{0.8em}
\hrule
\vspace{0.8em}

\subsubsection*{1. Hoạt động 1: Mở đầu / Khởi động (05 phút)}
\textbf{Chủ đề:} \textit{Trò chơi "Phân loại thần tốc": Ai là Cấp số cộng? Ai là Cấp số nhân?}
\begin{itemize}
    \item \textbf{a) Mục tiêu:} Tái hiện nhanh các dấu hiệu nhận biết một dãy số là cấp số cộng hay cấp số nhân; tạo không khí học tập sôi nổi và kích hoạt tư duy so sánh đối chiếu.
    \item \textbf{b) Nội dung:} Giáo viên chiếu lên màn hình 5 dãy số sau và yêu cầu học sinh giơ thẻ tín hiệu (hoặc trả lời nhanh trong 30 giây):
    \begin{itemize}
        \item (1) Dãy số: $3, 7, 11, 15, 19, \dots$
        \item (2) Dãy số: $3, 6, 12, 24, 48, \dots$
        \item (3) Dãy số: $5, 5, 5, 5, 5, \dots$
        \item (4) Dãy số: $2, -6, 18, -54, 162, \dots$
        \item (5) Dãy số: $1, 4, 9, 16, 25, \dots$
    \end{itemize}
    \item \textbf{c) Sản phẩm:} Câu trả lời của học sinh:
    \begin{itemize}
        \item (1) Cấp số cộng với $u_1 = 3$, công sai $d = 4$ (vì $7 - 3 = 11 - 7 = 4$).
        \item (2) Cấp số nhân với $u_1 = 3$, công bội $q = 2$ (vì $6/3 = 12/6 = 2$).
        \item (3) Vừa là cấp số cộng (với $d = 0$), vừa là cấp số nhân (với $q = 1$).
        \item (4) Cấp số nhân đan dấu với $u_1 = 2$, công bội $q = -3$ (vì $-6/2 = -3$).
        \item (5) Dãy số bình phương các số tự nhiên (không phải CSC, không phải CSN).
    \end{itemize}
    \item \textbf{d) Tổ chức thực hiện:} GV phổ biến luật chơi $\to$ HS quan sát, thảo luận cặp đôi chớp nhoáng $\to$ Đại diện xung phong trả lời và giải thích $\to$ GV chốt lại nguyên tắc cốt lõi: Phép cộng tạo ra CSC, phép nhân tạo ra CSN.
\end{itemize}

\subsubsection*{2. Hoạt động 2: Ôn tập hệ thống hóa kiến thức trọng tâm (12 phút)}

\paragraph{Mục 1: Bảng đối chiếu song song giữa Cấp số cộng và Cấp số nhân}

\begin{center}
\begin{tabularx}{\linewidth}{|l|X|X|}
\hline
\textbf{Đặc trưng so sánh} & \textbf{Cấp số cộng (CSC)} & \textbf{Cấp số nhân (CSN)} \\
\hline
\textbf{1. Định nghĩa truy hồi} & $u_{n+1} = u_n + d$ ($d$: công sai không đổi) & $u_{n+1} = u_n \cdot q$ ($q$: công bội không đổi) \\
\hline
\textbf{2. Cách xác định $d, q$} & $d = u_{n+1} - u_n$ (Phép trừ) & $q = \dfrac{u_{n+1}}{u_n}$ với $u_n \neq 0$ (Phép chia) \\
\hline
\textbf{3. Số hạng tổng quát} & $u_n = u_1 + (n - 1)d$ ($n \ge 2$) & $u_n = u_1 \cdot q^{n-1}$ ($n \ge 2$) \\
\hline
\textbf{4. Tính chất số hạng} & $u_k = \dfrac{u_{k-1} + u_{k+1}}{2}$ hay $u_{k-1} + u_{k+1} = 2u_k$ & $u_k^2 = u_{k-1} \cdot u_{k+1}$ hay $|u_k| = \sqrt{u_{k-1} u_{k+1}}$ \\
\hline
\textbf{5. Tổng $n$ số hạng đầu} & $S_n = \dfrac{n(u_1 + u_n)}{2} = \dfrac{n[2u_1 + (n - 1)d]}{2}$ & $S_n = \dfrac{u_1(1 - q^n)}{1 - q} = \dfrac{u_1(q^n - 1)}{q - 1}$ ($q \neq 1$) \\
\hline
\textbf{6. Dạng hàm số} & Tuyến tính bậc nhất: $f(n) = d \cdot n + (u_1 - d)$ & Hàm số mũ: $g(n) = \left(\dfrac{u_1}{q}\right) \cdot q^n$ \\
\hline
\end{tabularx}
\end{center}

\begin{scaffoldbox}[Giàn giáo sư phạm cho HS TB - Yếu: Mẹo tránh nhầm lẫn "kinh điển"]
    \begin{itemize}
        \item \textbf{Nhầm phép tính:} Cấp số cộng gắn với phép \textbf{cộng/trừ} ($d = u_{k+1} - u_k$), Cấp số nhân gắn với phép \textbf{nhân/chia} ($q = u_{k+1} / u_k$). Không bao giờ lấy $u_2 - u_1$ để tìm công bội của CSN!
        \item \textbf{Nhầm số mũ:} Trong công thức $u_n$ của CSN, số mũ của $q$ luôn là $(n - 1)$, không phải $n$. Ví dụ: $u_5 = u_1 \cdot q^4$.
        \item \textbf{Quên mẫu số của $S_n$:} Tổng CSC bắt buộc chia $2$: $S_n = \dfrac{n(u_1 + u_n)}{2}$. Tổng CSN có mẫu số là $(1 - q)$ hoặc $(q - 1)$ và chỉ áp dụng khi $q \neq 1$. Nếu $q = 1$ thì $S_n = n \cdot u_1$.
    \end{itemize}
\end{scaffoldbox}

\paragraph{Mục 2: Tích hợp Năng lực số (Mã 2.2NC1a) \& Năng lực AI (Mã 11.C3.1)}

\begin{pedabox}[Năng lực số (Mã 2.2NC1a): Chia sẻ tài nguyên học tập số]
    $\bullet$ Giáo viên yêu cầu đại diện các nhóm học sinh chụp ảnh sơ đồ tư duy hoặc tệp tóm tắt công thức đã chuẩn bị trước, gửi lên không gian lưu trữ Google Drive của lớp hoặc chia sẻ qua nhóm Zalo học tập.\\
    $\bullet$ Học sinh tải tệp tổng kết đối chiếu công thức về điện thoại cá nhân để làm cẩm nang tra cứu nhanh khi giải bài tập.
\end{pedabox}

\begin{aibox}[Tích hợp Năng lực AI (QĐ 2422/QĐ-BGDĐT --- Mã 11.C3.1): Kĩ thuật Prompt phân rã bài toán]
    $\bullet$ \textbf{Vấn đề thực tế:} Khi gặp bài toán hệ phương trình cấp số khó, học sinh thường sao chép toàn bộ đề bài và yêu cầu AI "Giải bài này cho em", dẫn đến việc tiếp thu thụ động và không hiểu bản chất.\\
    $\bullet$ \textbf{Kĩ thuật Prompt phân rã (Scaffolding Prompt):} Học sinh thực hành đặt prompt có cấu trúc như sau cho AI:\\
    \textit{"Bạn hãy đóng vai trò là một gia sư môn Toán lớp 11 tận tâm. Tôi đang giải bài toán: 'Tìm số hạng đầu $u_1$ và công sai $d$ của cấp số cộng thỏa mãn $u_2 + u_5 - u_3 = 10$ và $u_1 + u_6 = 17$'. Xin bạn ĐỪNG giải trực tiếp ra đáp án. Hãy đặt cho tôi 2 câu hỏi định hướng bước một để tôi tự quy đổi các số hạng $u_2, u_5, u_3, u_6$ về $u_1$ và $d$."}\\
    $\bullet$ Nhận diện phản hồi của AI: AI sẽ gợi ý công thức $u_n = u_1 + (n - 1)d$, từ đó học sinh tự tay viết phương trình và tìm ra đáp số.
\end{aibox}

\subsubsection*{3. Hoạt động 3: Luyện tập giải toán phân hóa (20 phút)}

\paragraph{Dạng 1: Xác định cấp số và tìm các yếu tố cơ bản (06 phút)}
\begin{itemize}
    \item \textbf{Bài toán 1:} Cho cấp số cộng $(u_n)$ có $u_1 = -3$ và $d = 5$.
    \begin{itemize}
        \item a) Tìm số hạng thứ $10$ của cấp số cộng.
        \item b) Số $497$ có thuộc cấp số cộng này không? Nếu có, là số hạng thứ mấy?
        \item c) Tính tổng của $30$ số hạng đầu tiên ($S_{30}$).
    \end{itemize}
    \item \textit{Lời giải chi tiết:}\\
    a) Áp dụng công thức số hạng tổng quát: $u_{10} = u_1 + 9d = -3 + 9 \cdot 5 = 42$.\\
    b) Giả sử $u_n = 497 \iff -3 + (n - 1) \cdot 5 = 497 \iff 5(n - 1) = 500 \iff n - 1 = 100 \iff n = 101 \in \mathbb{N}^*$.\\
    Vậy số $497$ là số hạng thứ $101$ của cấp số cộng ($u_{101} = 497$).\\
    c) Áp dụng công thức tính tổng $S_n$:
    $$S_{30} = \dfrac{30 \cdot [2 \cdot (-3) + (30 - 1) \cdot 5]}{2} = 15 \cdot (-6 + 29 \cdot 5) = 15 \cdot 139 = 2085$$
\end{itemize}

\paragraph{Dạng 2: Thiết lập và giải hệ phương trình tìm cấp số (08 phút)}
\begin{itemize}
    \item \textbf{Bài toán 2:} Tìm số hạng đầu $u_1$ và công bội $q$ của cấp số nhân $(u_n)$ biết:
    $$\begin{cases} u_1 + u_3 = 10 \\ u_4 + u_6 = 80 \end{cases}$$
    \item \textit{Lời giải chi tiết:}\\
    Biểu diễn các số hạng theo $u_1$ và $q$ ($u_k = u_1 \cdot q^{k-1}$):
    $$\begin{cases} u_1 + u_1 q^2 = 10 \\ u_1 q^3 + u_1 q^5 = 80 \end{cases} \iff \begin{cases} u_1(1 + q^2) = 10 \quad (1) \\ u_1 q^3(1 + q^2) = 80 \quad (2) \end{cases}$$
    Vì $1 + q^2 > 0$ với mọi $q$, từ $(1)$ suy ra $u_1 \neq 0$.\\
    Lấy phương trình $(2)$ chia cho phương trình $(1)$ vế theo vế:
    $$\dfrac{u_1 q^3 (1 + q^2)}{u_1 (1 + q^2)} = \dfrac{80}{10} \iff q^3 = 8 \iff q = 2$$
    Thay $q = 2$ vào phương trình $(1)$:
    $$u_1 (1 + 2^2) = 10 \iff 5u_1 = 10 \iff u_1 = 2$$
    Vậy cấp số nhân có số hạng đầu $u_1 = 2$ và công bội $q = 2$.
\end{itemize}

\paragraph{Dạng 3: Bài toán kết hợp giữa Cấp số cộng và Cấp số nhân (06 phút)}
\begin{itemize}
    \item \textbf{Bài toán 3:} Cho ba số $a, b, c$ theo thứ tự lập thành một cấp số cộng có tổng bằng $15$. Nếu ta lần lượt cộng thêm $1, 4, 19$ vào ba số đó thì được ba số mới theo thứ tự lập thành một cấp số nhân. Tìm ba số $a, b, c$.
    \item \textit{Lời giải chi tiết:}\\
    Do $a, b, c$ lập thành cấp số cộng nên ta có: $a + c = 2b$.\\
    Mặt khác, tổng ba số bằng $15$:
    $$(a + c) + b = 15 \iff 2b + b = 15 \iff 3b = 15 \iff b = 5$$
    Gọi $d$ là công sai của cấp số cộng, ta có: $a = 5 - d$ và $c = 5 + d$.\\
    Khi cộng thêm $1, 4, 19$ vào ba số, ta được ba số mới:
    \begin{itemize}
        \item $x = a + 1 = (5 - d) + 1 = 6 - d$
        \item $y = b + 4 = 5 + 4 = 9$
        \item $z = c + 19 = (5 + d) + 19 = 24 + d$
    \end{itemize}
    Ba số $x, y, z$ lập thành cấp số nhân khi và chỉ khi:
    $$y^2 = x \cdot z \iff 9^2 = (6 - d)(24 + d)$$
    $$\iff 81 = 144 + 6d - 24d - d^2 \iff d^2 + 18d - 63 = 0$$
    Giải phương trình bậc hai trên máy tính cầm tay (\texttt{MENU} 9-2-2):
    $$\Delta' = 9^2 - 1 \cdot (-63) = 81 + 63 = 144 = 12^2 \implies \begin{bmatrix} d = -9 + 12 = 3 \\ d = -9 - 12 = -21 \end{bmatrix}$$
    \begin{itemize}
        \item Với $d = 3$: Ta có $a = 5 - 3 = 2$, $b = 5$, $c = 5 + 3 = 8$. (Ba số mới là $3, 9, 27$ là CSN với $q = 3$).
        \item Với $d = -21$: Ta có $a = 5 - (-21) = 26$, $b = 5$, $c = 5 - 21 = -16$. (Ba số mới là $27, 9, 3$ là CSN với $q = \dfrac{1}{3}$).
    \end{itemize}
    Kết luận: Có hai bộ ba số thỏa mãn đề bài là $(2; 5; 8)$ hoặc $(26; 5; -16)$.
\end{itemize}

\subsubsection*{4. Hoạt động 4: Vận dụng thực tiễn (08 phút)}
\textbf{Chủ đề:} \textit{So sánh bài toán tích lũy tài chính: Mô hình Tuyến tính (CSC) vs Mô hình Lãi kép (CSN)}
\begin{itemize}
    \item \textbf{a) Mục tiêu:} Vận dụng kiến thức tổng hợp của cấp số cộng và cấp số nhân vào việc ra quyết định tài chính cá nhân; khắc sâu sức mạnh tăng trưởng hàm mũ của lãi kép trong dài hạn.
    \item \textbf{b) Nội dung:} Anh Bình và anh Cường cùng có số vốn ban đầu là $100$ triệu đồng để đầu tư trong thời hạn $10$ năm:
    \begin{itemize}
        \item \textbf{Phương án A (Anh Bình):} Đầu tư sản xuất với lợi nhuận cố định mỗi năm đều đặn là $12$ triệu đồng (lợi nhuận không tái đầu tư). Sau mỗi năm, tổng tài sản tăng thêm đúng $12$ triệu đồng (mô hình Cấp số cộng).
        \item \textbf{Phương án B (Anh Cường):} Gửi tiết kiệm ngân hàng theo hình thức lãi kép với lãi suất ổn định $8\%$/năm, toàn bộ tiền lãi được tự động nhập gốc để tính lãi cho các năm tiếp theo (mô hình Cấp số nhân).
    \end{itemize}
    Hỏi sau $10$ năm, tổng tài sản của ai lớn hơn và chênh lệch nhau bao nhiêu tiền?
    \item \textbf{c) Sản phẩm:}\\
    \textit{Lời giải:}\\
    - Đối với anh Bình: Tổng tài sản sau mỗi năm lập thành cấp số cộng với $u_0 = 100$ triệu, công sai $d = 12$ triệu.\\
    Sau $10$ năm, số tiền của anh Bình là:
    $$T_A = 100 + 10 \cdot 12 = 100 + 120 = 220\text{ (triệu đồng)}$$
    - Đối với anh Cường: Số tiền sau mỗi năm lập thành cấp số nhân với số vốn ban đầu $V_0 = 100$ triệu, công bội $q = 1 + 0{,}08 = 1{,}08$.\\
    Sau $10$ năm, số tiền của anh Cường là:
    $$T_B = 100 \cdot (1{,}08)^{10} \approx 100 \cdot 2{,}158925 = 215{,}89\text{ (triệu đồng)}$$
    - Nhận xét và so sánh:
    \begin{itemize}
        \item Sau $10$ năm: Anh Bình có $220$ triệu đồng, anh Cường có $215{,}89$ triệu đồng. Anh Bình nhiều hơn anh Cường khoảng $4{,}11$ triệu đồng.
        \item Tuy nhiên, nếu xét sau $15$ năm:
        Anh Bình có: $100 + 15 \cdot 12 = 280$ triệu đồng.\\
        Anh Cường có: $100 \cdot (1{,}08)^{15} \approx 317{,}22$ triệu đồng $\implies$ Anh Cường vượt trội hơn anh Bình tới $37{,}22$ triệu đồng!
    \end{itemize}
    $\implies$ \textbf{Bài học kinh tế:} Cấp số cộng có ưu thế trong ngắn hạn, nhưng cấp số nhân (lãi kép) sẽ bùng nổ vượt trội hoàn toàn trong dài hạn.
    \item \textbf{d) Tổ chức thực hiện:} GV giao bài toán thực tế $\to$ HS phân tích mô hình đại số tương ứng $\to$ Bấm máy tính so sánh $\to$ GV tổng kết ý nghĩa thực tiễn của Chương II.
\end{itemize}

\newpage

\section*{IV. PHỤ LỤC HỌC LIỆU BỔ TRỢ}

\subsection*{Phụ lục 1: Phiếu học tập số 1 --- Bảng đối chiếu tự hoàn thiện CSC và CSN}
\noindent\textbf{Họ và tên:} .................................................................... \textbf{Lớp:} 11A....... \textbf{Nhóm:} ..........

\begin{pedabox}[Nhiệm vụ: Điền công thức tương ứng vào các ô còn trống]
\begin{center}
\begin{tabularx}{0.95\linewidth}{|p{3.8cm}|X|X|}
\hline
\textbf{Nội dung} & \textbf{Cấp số cộng} & \textbf{Cấp số nhân} \\
\hline
Công thức số hạng tổng quát $u_n$ & $u_n = $ .............................. & $u_n = $ .............................. \\
\hline
Quan hệ giữa 3 số hạng liên tiếp & ........................................... & ........................................... \\
\hline
Công thức tính tổng $n$ số hạng đầu $S_n$ & $S_n = $ .............................. & $S_n = $ .............................. \\
\hline
Điều kiện áp dụng & Với mọi $d$ & Bắt buộc $q \neq $ .............. \\
\hline
\end{tabularx}
\end{center}
\end{pedabox}

\subsection*{Phụ lục 2: Phiếu học tập số 2 --- Bài tập tự luyện phân hóa}
\noindent\textbf{Họ và tên:} .................................................................... \textbf{Lớp:} 11A....... \textbf{Nhóm:} ..........

\begin{pedabox}[Bộ bài tập rèn luyện]
    $\bullet$ \textbf{Bài 1 (Cơ bản):} Cho cấp số cộng $(u_n)$ có $u_1 = 4, d = -3$. Tính $u_{15}$ và $S_{20}$.\\
    $\bullet$ \textbf{Bài 2 (Trung bình):} Cho cấp số nhân $(u_n)$ có $u_2 = 6$ và $u_5 = 162$. Tìm $u_1$ và công bội $q$.\\
    $\bullet$ \textbf{Bài 3 (Vận dụng):} Tìm ba số hạng liên tiếp của một cấp số cộng, biết tổng của chúng bằng $21$ và tổng các bình phương của chúng bằng $197$.\\
    $\bullet$ \textbf{Bài 4 (AI Prompt):} Hãy viết một câu prompt ngắn yêu cầu trợ lý AI giải thích vì sao một dãy số vừa có thể là cấp số cộng, vừa có thể là cấp số nhân.
\end{pedabox}

\subsection*{Phụ lục 3: Bảng Rubric đánh giá năng lực tích hợp trong Bài tập cuối chương II}

\begin{center}
\begin{tabularx}{\linewidth}{|p{3.2cm}|X|X|X|X|}
\hline
\textbf{Tiêu chí} & \textbf{Mức 1 (Chưa đạt)} & \textbf{Mức 2 (Đạt)} & \textbf{Mức 3 (Khá)} & \textbf{Mức 4 (Tốt)} \\
\hline
\textbf{1. Hệ thống hóa kiến thức Chương II} & Nhầm lẫn công thức giữa CSC và CSN; chưa phân biệt được $d$ và $q$. & Nhớ và áp dụng đúng các công thức $u_n, S_n$ cho từng cấp số riêng lẻ. & Phân biệt chính xác tính chất đặc trưng của CSC và CSN; giải đúng hệ phương trình. & Vận dụng linh hoạt giải quyết trọn vẹn các bài toán kết hợp đan xen giữa CSC và CSN. \\
\hline
\textbf{2. Ứng dụng số (Mã 2.2NC1a)} & Chưa biết cách chia sẻ tệp tài liệu số qua mạng. & Tải và lưu trữ được tài liệu ôn tập từ Drive/Zalo nhóm lớp. & Chủ động chia sẻ sơ đồ tóm tắt công thức cho bạn học qua nhóm số. & Khai thác hiệu quả môi trường số để thảo luận, giải đáp bài tập trực tuyến. \\
\hline
\textbf{3. Năng lực AI (Mã 11.C3.1)} & Chỉ sao chép nguyên văn đề bài yêu cầu AI giải hộ. & Bước đầu biết đặt prompt yêu cầu AI giải thích công thức. & Thực hành thành thạo prompt phân rã bài toán thành các câu hỏi gợi ý định hướng. & Đánh giá và kiểm chứng độc lập lời giải của AI, phát hiện các sai số logic nếu có. \\
\hline
\end{tabularx}
\end{center}

\end{document}
